\begin{claim}{D.68}{Proved}{The exact dihedral quotient of a quartic torus action}
Let $a$ have order four in the full once-punctured-torus
loop-outer group. Choose an actual order-four representative
$\alpha$ as in D.46, and write
\[
 F=\widehat F(X,Y),\qquad c=[X,Y],\qquad
 E=F\rtimes_\alpha\langle s\rangle .
\]
There is exactly one nonzero character $\chi:F\to C_2$
fixed by $a$. Extend it to $\widetilde\chi:E\to C_2$
by $\widetilde\chi(s)=0$. Every cusp-normalizer element
$v$ with quotient label $a$ has
\[
                   \widetilde\chi(v)=1.                    \tag{D.68a}
\]

There is an explicitly specified discrete free basis $x,y$
with $[x,y]=c$, $\chi(x)=1$, and $\chi(y)=0$.
For this basis the closed normal subgroup
\[
                K=\langle\!\langle x^2,(xy)^2
                           \rangle\!\rangle_F              \tag{D.68b}
\]
is $\alpha$-invariant. The actual quotient and its inverse
presentation maps identify
\[
 \begin{split}
 F/K&\cong\mathcal D=\widehat\Z\rtimes C_2
       =\langle r,t\mid t^2=1, trt=r^{-1}\rangle,\\
 x&\longmapsto t,\qquad y\longmapsto r,
                 \qquad c\longmapsto r^{-2}.
 \end{split}                                               \tag{D.68c}
\]
Here $r$ generates the full $\widehat\Z$ rotation subgroup.
The induced automorphism is exactly
\[
 \overline\alpha(r)=r^{-\lambda},\qquad
 \overline\alpha(t)=r^j t,\qquad
 j_2\in\Z_2^\times,\qquad (1-\lambda)j=0.                \tag{D.68d}
\]
It has order two. For every odd prime $q$ the further
specified quotient $\mathcal D\to D_{2q}$ sends $r$ to
a rotation of order $q$ and $t$ to a reflection.
On that quotient the induced action is the identity
exactly at the coordinates $\lambda_q=-1$; at
$\lambda_q=1$ it is an involution inverting the rotation.
In particular each remaining scalar quartic coordinate
has an actual nonabelian quotient $F\twoheadrightarrow D_{2q}$
fixed pointwise by this representative.

The proof also gives the full rank-three homology matrix
of the actual index-two cover, with all cusp labels retained.
These conclusions constrain every actual quartic action;
they do not construct a negative multiplier coordinate
or settle its remaining exclusion or conjugacy problem.
\end{claim}
\begin{proof}
\emph{The specified mod-two character and a cusp-preserving basis.}
By D.55 the positive square of $a$ acts as $-I$ on
$H_1(F,\Z_p)$ for every prime. Hence
\[
 A_p(a)^2=-I,\qquad \lambda^2=1,\qquad \lambda_2=1.       \tag{D.68e}
\]
The last assertion also follows directly from the
characteristic polynomial $Z^2+1$ over $\Q_2$ and
the determinant interpretation of the cusp multiplier.
The matrix $A_2(a)$ cannot be identity modulo two:
otherwise $A_2(a)=I+2B$ would have square identity
modulo four, contradicting its square $-I$.
Its reduction is therefore a nonidentity unipotent
matrix over $\F_2$. Its fixed dual line has exactly
one nonzero vector. This is the stated character $\chi$.
Conjugation on $F$ does not change a character, so
this character is defined by the outer action as well.

Here are the basis maps and their actual inverses,
according to the values of $\chi$ on $X,Y$:
\[
\begin{array}{c|cc|cc}
 (\chi(X),\chi(Y))&x&y&X&Y\\ \hline
 (1,0)&X&Y&x&y\\
 (0,1)&XYX^{-1}&X^{-1}&y^{-1}&yxy^{-1}\\
 (1,1)&X&YX^{-1}&x&yx.
\end{array}                                                \tag{D.68f}
\]
Free reduction verifies both composites and $[x,y]=[X,Y]$
in every row. The same substitutions give mutually
inverse maps on full free profinite groups. No peripheral
conjugation or change of multiplier has been introduced.

Put $N=\ker\chi$. Its actual discrete double cover of
the punctured torus has type $(1,2)$. Schreier rewriting
with transversal $1,x$ gives its full free profinite basis
\[
 u=x^2,\qquad B=y,\qquad C=xyx^{-1},\qquad c=CB^{-1}.
                                                               \tag{D.68g}
\]
The second based peripheral generator is $xcx^{-1}$;
its homology class is the negative of that of $c$.
Write $M=H_1(N,\widehat\Z)$ in the ordered basis
$[u],[B],[C]$ and put $b=[C]-[B]$.
The actual deck action $J$, induced by conjugation
by $x$, is
\[
 J[u]=[u],\qquad J[B]=[C],\qquad J[C]=[B].                \tag{D.68h}
\]
Before abelianization the last identity is
$xCx^{-1}=uBu^{-1}$. Homology of the actual inclusion
$N\hookrightarrow F$ is the map
\[
 I:M\longrightarrow H_1(F,\widehat\Z),\qquad
 I=\begin{pmatrix}2&0&0\\0&1&1\end{pmatrix},\qquad
 \ker I=\widehat\Z b.                                    \tag{D.68i}
\]
It factors through filling both cusps of the covered
torus and then the index-two closed-torus lattice
inclusion. Its image is precisely that lattice.
Thus after
tensoring each primary coordinate with $\Q_p$, the
$+1$ eigenspace of $J$ is the handle space and the
$-1$ eigenspace is precisely $\Q_p b$.

\emph{The cusp root has the nonzero character label.}
Choose $v$ above $a$ in the normalizer of the actual
cusp subgroup. The cusp relations, including their
primary exponent, are
\[
 vcv^{-1}=c^\lambda,\qquad v^4=c^\kappa,\qquad
                         \kappa_2\in\Z_2^\times.          \tag{D.68j}
\]
The unit assertion is D.55 applied to $v^2$, which
lies over the positive involution $a^2$.
The invariant character extends to $E$ by sending $s$
to zero. Since $\chi$ is invariant, $N$ is normal in
$E$. Let $T$ be conjugation by $v$ on $M$.
The element $vxv^{-1}$ has the same nonzero character
as $x$, so its action on $M$ is also $J$.
Consequently $T$ commutes with $J$.

On the boundary line $T$ is multiplication by $\lambda$.
On the filled handle quotient it induces the actual
matrix $A(a)$. Equations (D.68e) and (D.68i), first
over $\Q_p$ and then in the torsion-free integral
module for each $p$, prove
\[
                   T^2=-J.                                \tag{D.68k}
\]
At the two-primary coordinate $Tb=b$, and therefore
\[
 \mathcal N_4(T)=(I+T)(I+T^2)
           =(I+T)(I-J)=2(I-J),\qquad
 \mathcal N_4(T)M_2=2\Z_2b.                               \tag{D.68l}
\]
If $\widetilde\chi(v)$ were zero, write $v=ns$ with
$n\in N$. Conjugation by $v$ and by $s$ then agree
on $M$. The actual identity
$(ns)^4=n\alpha(n)\alpha^2(n)\alpha^3(n)$ would give
$\kappa_2b=\mathcal N_4(T)[n]$, contrary to
(D.68l) and the unit exponent in (D.68j).
This proves (D.68a).

Let $S$ denote the action of $\alpha$ on $M$.
Writing $v=f_0s$, the proved label means $\chi(f_0)=1$,
so conjugation by $f_0$ on $M$ is exactly $J$.
Thus the actual actions satisfy
\[
            T=JS,\qquad S=JT,\qquad
            S^2=-J,\qquad Sb=-\lambda b.                  \tag{D.68m}
\]
In particular the primary norm of $S$ at two is zero.

The entire matrix is explicit. Write the base homology
matrix in the basis $[x],[y]$ as
$A(a)=\left(\begin{smallmatrix}a_{11}&a_{12}\\
a_{21}&a_{22}\end{smallmatrix}\right)$. Then
\[
 T=\begin{pmatrix}
 a_{11}&a_{12}/2&a_{12}/2\\
 a_{21}&(a_{22}+\lambda)/2&(a_{22}-\lambda)/2\\
 a_{21}&(a_{22}-\lambda)/2&(a_{22}+\lambda)/2
 \end{pmatrix},\qquad S=JT.                               \tag{D.68n}
\]
Every half coefficient belongs to $\widehat\Z$:
the invariant mod-two character makes $a_{12}$ even,
and $a_{22}$ and $\lambda$ are odd. To verify the
matrix without choosing an integral splitting, use
$TJ=JT$, $Tb=\lambda b$, and the actual equality
$IT=A(a)I$. These equations determine the displayed
entries uniquely.

\emph{The full dihedral quotient and both presentation maps.}
Define the actual character
\[
 \ell:N\longrightarrow\widehat\Z,\qquad
       \ell(u)=0,\quad\ell(B)=1,\quad\ell(C)=-1.
                                                               \tag{D.68o}
\]
It is surjective, and $\ell\circ\operatorname{Inn}(x)=-\ell$.
Thus its kernel is normal in $F$. Equation (D.68n)
gives, as an equality of characters of $N$,
\[
                         \ell\alpha=-\lambda\ell.         \tag{D.68p}
\]
Since $-\lambda$ is a unit, $\ker\ell$ is
$\alpha$-invariant as a subgroup, in both directions.

The homomorphism $F\to\mathcal D$ sending $x$ to $t$
and $y$ to $r$ restricts on $N$ to $n\mapsto r^{\ell(n)}$.
It is onto, and its kernel is $\ker\ell$. Indeed an
element outside $N$ maps to the nontrivial reflection
coset, while the rotation parameter detects every
element of $N/\ker\ell$.
This is also precisely the quotient (D.68b): in that
quotient $x^2=1$ and $(xy)^2=1$, so
$xyx=y^{-1}$. Conversely the images $x=t,y=r$ in
$\mathcal D$ satisfy these two relators.
The inverse homomorphism sends $r$ to $yK$ and $t$
to $xK$. Both composites fix the named generators.
The rotation subgroup is full $\widehat\Z$, also
directly witnessed by the surjective character $\ell$.
This proves all assertions about the actual kernel
and maps in (D.68b)--(D.68c).

Equation (D.68p) specifies the image of $r$ under the
induced automorphism. Preservation of $\chi$ specifies
the reflection coset of the image of $t$. Hence there
is a unique $j\in\widehat\Z$ with the two formulas
in (D.68d). Their square sends $r$ to $r$ and $t$ to
$r^{(1-\lambda)j}t$. Since $\alpha^4=1$, its fourth
power gives $2(1-\lambda)j=0$ in $\widehat\Z$.
Multiplication by two is injective, proving the last
equation of (D.68d) and $\overline\alpha^2=1$.

For the retained cusp label one can compute $j$ more
precisely. The quotient extends to the actual map
\[
 E\longrightarrow\mathcal D\rtimes_{\overline\alpha}C_4,
 \qquad f\longmapsto\pi(f),\quad s\longmapsto s.
                                                               \tag{D.68q}
\]
Write $\pi(f_0)=r^m t$, using $\chi(f_0)=1$.
In this quotient the exact square is
\[
 \pi(v)^2=r^{(1+\lambda)m-j}s^2,\qquad
 \pi(v)^4=r^{2((1+\lambda)m-j)}=r^{-2\kappa}.
\]
Injectivity of multiplication by two in the rotation
parameter gives
\[
                         j=(1+\lambda)m+\kappa.            \tag{D.68r}
\]
At two, this reads $j_2=2m_2+\kappa_2$, so $j_2$ is
odd. In particular the induced action is nonidentity
and has exact order two. The formula retains the
original exponent and does not rescale the cusp.

Finally reduce the full rotation parameter modulo an
odd prime $q$. If $\lambda_q=-1$, equation
$(1-\lambda)j=0$ gives $j_q=0$ already in $\Z_q$.
Both displayed generators of $D_{2q}$ are fixed.
If $\lambda_q=1$, the rotation is inverted and is
nontrivial because $q$ is odd, so the induced action
has exact order two. The cusp has image $r^{-2}$,
which still generates the order-$q$ rotation subgroup.
Thus these are actual peripheral finite quotients,
not quotients inferred only from a dimension count.

The executed certificate
\path{code/part_c/torus_quartic_dihedral_d68.py}
checks the discrete basis maps, cusp and deck words,
the double-cover matrices and integral norm identity,
the standard quartic representative, and the exact
dihedral formulas at finite primary coordinates.
It does not assert that a permitted coordinate
model lifts to a full torus automorphism.
\end{proof}
